\chapter{}
\label{chap:3}

\section{Умова}
Нехай $f \in B_{n}$, $\AI(f) = d$. Довести, що:
\begin{equation*}
    \sum\limits_{i=0}^{d-1} \binom{n}{i} \leq \Vert f \Vert \leq \sum\limits_{i=0}^{n-d} \binom{n}{i}
\end{equation*}

\section{Розв'язання}
Мені повезло, що я ходив на практики і робив їх конспект)) Тож я просто переписав доведення, яке ми зробили протягом 
четвертої практики в \LaTeX{}. Тут доведення складається з двох частин: спочатку ми доводили ліву нерівність (оцінку 
знизу), а потім оцінку зверху.

\subsection*{Частина 1}

Розглянемо матрицю $M$, рядки якої пронумеровані векторами $x \in V_{n}$, для яких функція набуває значення 
один ($f(x) = 1$), а стовпці матриці відповідають усім можливим мономам степеня від $0$ до $d-1$ включно. На 
перетині рядка $x$ та стовпця -- значення цього монома в точці $x$. Підрахуємо розмірність цієї матриці. 
Кількість рядків дорівнює кількості одиниць у векторі значень функції, тобто вазі Геммінга функції $f$:
\begin{equation*}
    \text{К-ть рядків} = \Vert f \Vert
\end{equation*}
А кількість стовпців дорівнює відповідно загальній кількості мономів степеня $\leq d-1$ від $n$ змінних. Кількість 
мономів степеня $i$ дорівнює кількості комбінацій $\binom{n}{i}$, тому:
\begin{equation*}
    \text{К-ть стовпців} = \sum\limits_{i=0}^{d-1} \binom{n}{i}
\end{equation*}
Ми доводили від супротивного. Нехай оцінка знизу не виконується, тобто кількість стовпців є строго більшою за кількість 
рядків:
\begin{equation*}
    \sum\limits_{i=0}^{d-1} \binom{n}{i} > \Vert f \Vert
\end{equation*}
З курсу лінійки відомо, що якщо в однорідній системі лінійних рівнянь кількість невідомих (стовпців) більша за кількість 
рівнянь (рядків), то така система гарантовано має ненульовий розв'язок. Отже, матричне рівняння $Mz = 0$ має ненульовий 
розв'язок $z$. Цей вектор-стовпець $z$ визначає коефіцієнти деякого ненульової функції $g \in B_{n}$, і при цьому 
$g \in \Ann(f)$. З побудови матриці $M$ випливає:
\begin{enumerate}
    \item Максимальний степінь цього полінома обмежений вибраними стовпцями: $\deg(g) \leq d-1$.
    \item Оскільки $Mz = 0$, то для всіх $x$, де $f(x)=1$, виконується $g(x) = 0$. Тобто добуток тотожно дорівнює нулю: 
        $f \cdot g = 0$, а отже $g$ є анулятором функції $f$ ($g \in Ann(f)$).
\end{enumerate}
А якщо у функції $f$ існує анулятор $g$ степеня $\leq d-1$, то за визначенням алгебраїчної імунності:
\begin{equation*}
    AI(f) \leq d-1
\end{equation*}
Але за умовою задачі $AI(f) = d$. Ми отримали суперечність. Отже, припущення було хибним, кількість рядків має 
бути не меншою за кількість стовпців.
\begin{equation*}
    \Vert f \Vert \geq \sum\limits_{i=0}^{d-1} \binom{n}{i}.
\end{equation*}

\subsection*{Частина 2}
Тут ми користувалися функцією $f \oplus 1$, і робили аналогічні кроки. Алгебраїчна імунність функції та її 
інверсії ($\oplus 1$) завжди збігаються, оскільки будь-який анулятор для $f$ є анулятором для $f \oplus 1$ (і навпаки). 
Отже, $AI(f \oplus 1) = d$. Застосуємо до функції $f \oplus 1$ щойно доведену оцінку знизу:
\begin{equation*}
    \Vert f \oplus 1 \Vert \geq \sum\limits_{i=0}^{d-1} \binom{n}{i}
\end{equation*}
Вага Геммінга для такої функції обчислюється як різниця загальної кількості векторів простору мінус вага початкової 
функції: $\Vert f \oplus 1 \Vert = 2^{n} - \Vert f \Vert$. Підставимо це у нерівність:
\begin{equation*}
    2^{n} - \Vert f \Vert \geq \sum\limits_{i=0}^{d-1} \binom{n}{i}
\end{equation*}
Змінимо місцями доданки:
\begin{equation*}
    \Vert f \Vert \leq 2^{n} - \sum\limits_{i=0}^{d-1} \binom{n}{i}
\end{equation*}
І пригадаємо, що сума всіх комбінацій дорівнює $2^{n}$, тобто:
\begin{equation*}
    \sum\limits_{i=0}^{n} \binom{n}{i} = 2^{n}
\end{equation*}
Підставимо і віднімемо:
\begin{equation*}
    \Vert f \Vert \leq \sum\limits_{i=0}^{n} \binom{n}{i} - \sum\limits_{i=0}^{d-1} \binom{n}{i} = \sum\limits_{i=d}^{n} \binom{n}{i}
\end{equation*}
За властивістю біноміальних коефіцієнтів: $\binom{n}{i} = \binom{n}{n-i}$. Тому:
\begin{equation*}
    \Vert f \Vert \leq \sum\limits_{i=0}^{n-d} \binom{n}{i}
\end{equation*}
Фінально отримали те, що й мали:
\begin{equation*}
    \sum\limits_{i=0}^{d-1} \binom{n}{i} \leq \Vert f \Vert \leq \sum\limits_{i=0}^{n-d} \binom{n}{i}
\end{equation*}
